% !TEX root = ../main.tex

\subsection{Headline result and comparison with prior work}

The final blended ranker reaches NDCG@5 = 0.41521 on a locked 19{,}980-search holdout and 0.41706 on the Kaggle public leaderboard. The result is broadly comparable to the Bellucci~\cite{bellucci2021dmt} reference of 0.41726 on the same in-class evaluation, despite three notable differences in approach: OOF rather than LOO target encoding, the addition of an XGBoost blend partner with a Wang-Kalousis~\cite{wang2014expedia} style of within-search relativisations, and the omission of the Bellucci per-property positional proxy from \texttt{random\_bool} = 1 rows. The interpretation we take from this is that the bulk of the public-leaderboard signal on this dataset is captured by a careful but minimal set of within-search relativisations and \texttt{prop\_id}-keyed priors; adding more features beyond block E or replacing OOF by LOO is unlikely to move the headline by more than 0.003 to 0.005 NDCG@5 on this evaluation.

\subsection{A note on the choice of NDCG@5}

The course rubric fixes NDCG@5 with the IDCG = 0 convention as the evaluation metric, which is a choice with concrete consequences for how we read the headline number. NDCG@5 was introduced by J\"arvelin and Kek\"al\"ainen~\cite{jarvelin2002ndcg} and uses a logarithmic positional discount that mirrors a user's decreasing attention down a results list. Three alternatives would have produced a meaningfully different ranking of our predictors. Precision@5 ignores the order within the top five entirely, which would have flattened the small but real differences between our LightGBM seeds, since the top-five sets agree on most searches even when the within-top-five order does not. Mean Reciprocal Rank rewards only the top-ranked positive, which would have weighted booked hotels much more heavily than clicked ones since 95\% of searches with positives have at most one click and one booking. Expected Reciprocal Rank, introduced by Chapelle et al.~\cite{chapelle2009expected}, models the user as stopping after a satisfied click, which on this dataset would have penalised the medium-popularity tertile more than NDCG@5 does and would have made the fairness rerank look stronger. The IDCG = 0 convention itself depresses our reported NDCG@5 by roughly 0.04 in absolute terms relative to the "skip empty searches" convention. The convention is appropriate for an in-class leaderboard, because it treats a search with no booked or clicked hotel as a genuine failure of the ranker rather than as missing data, but it is worth knowing that the headline can be inflated or deflated by 0.03 to 0.05 NDCG@5 by the convention alone.

\subsection{Limitations}

Four limitations of the present pipeline are worth flagging. First, we did not implement Bellucci's per-property positional proxy from \texttt{random\_bool} = 1 rows. The EDA distribution of position on shuffled pages is not uniform, so the assumption underlying the proxy does not strictly hold; a pilot of the proxy yielded a gain below noise. A more careful version, conditional on page size, might extract some signal. Second, the blend weight grid is coarse (eight discrete weights). A finer grid would likely move $w_{\text{xgb}}$ by 0.01 to 0.05, but the val plateau between 0.5 and 0.7 is flat enough that the headline NDCG would not move materially. Third, the fairness rerank uses a single regularisation strength $\lambda$ and a uniform $-\log P_t$ boost across all items of a tertile. A more granular formulation, with a per-search rather than a per-tertile re-balance, would presumably extract more medium-tertile lift at the same global cost; we did not implement it because the gain on top of the already balanced Optuna model is bounded. Fourth, the local-to-Kaggle shrinkage analysis is based on a single submission. A leave-one-out cross-validation of the shrinkage over multiple submissions would give a defensible confidence interval rather than a point.

\subsection{Future work}

Three directions stand out. First, replacing OOF target encoding by LOO at the \texttt{prop\_id} level, which Bellucci~\cite{bellucci2021dmt} reported as their largest single feature contributor; we used OOF for tractability, but the gap to LOO has not been measured carefully on this dataset. Second, a learned-position-bias correction in the loss function rather than a sample-weighting heuristic. The sample-weighting experiment of Section~\ref{sec:models} hurt NDCG, but a principled propensity-weighted lambda gradient, in the style of unbiased learning-to-rank, might recover the random-page gap of 0.086 that the current model leaves on the table. Third, a per-user or per-session feature space, which the current dataset only partially exposes through \texttt{visitor\_hist\_*}. Worst-case characterisation in Section~\ref{sec:eval} indicates that 40\% of holdout searches score exactly NDCG@5 = 0, and the bottleneck is genuinely the absence of a per-user prior.

\subsection{Reproducibility}

All code is organised as six numbered notebooks in \texttt{notebooks/}. Running them in order from \texttt{01\_eda.ipynb} to \texttt{06\_final\_pipeline\_and\_submission.ipynb} reproduces every number in this report from the raw CSVs in \texttt{data/raw/}. The locked split is stored in \texttt{data/processed/split\_seed42.json}; the Optuna best parameters are in \texttt{optuna\_lgbm\_best.json} and \texttt{optuna\_xgb\_best.json}; the realised Kaggle score is in \texttt{kaggle\_public\_score.json}. Every figure and table in this report cites a specific file on disk so a reviewer can verify the numbers without re-running the pipeline. Memory caps are explicit (\texttt{n\_jobs = 4} throughout, never $-1$, with \texttt{gc.collect()} between model fits) so the pipeline runs end-to-end on a 64 GB laptop in under three hours.
